
\begin{center}
    {\Large\itshape Alban pendant la guerre de Cent Ans}
\end{center}

\vspace{1cm}

À quelque temps de là, Albi est informé que des bandes inconnues
occupent Villeneuve du Tarn. En mars 1361, les consuls d'Alban et
ceux de Villefranche avisent Albi à se tenir sur ses gardes.
Deux mois après, les consuls d'Albi envoient des émissaires
s'informer auprès des consuls d'Alban.

\medskip

En 1362, châteaux et bourgs des environs tombent entre les mains
de tous ces pillards. Massals, Miolles, Saint Jean de Janes,
Massuys, Curvalle, Trébas, Gaycre, Aygou, Padiès, Rosières,
Tanus\ldots

\medskip

En septembre 1368, Alban tombe aux mains des compagnies de
Provence, puis c'est le tour de Villefranche, et le 29 septembre
de cette même année, elles entrent dans les faubourgs d'Albi.

\medskip

Un observateur était envoyé, en juillet 1370, pour
« espiar »\footnote{Épier.} et s'informer sur les troupes qui venaient de
franchir le Tarn au bac de Trébas.

\medskip

Un an après, en 1371, au mois de mars, le trésorier du comte
d'Armagnac est à Alban. Un sergent nommé « Lyéric Davy »
envoyé par les consuls d'Albi était venu l'y rejoindre,
pour réclamer une paire de draps qui avaient été prêtés
et n'avaient pas été rendus.
